\section{Simulation methods}
\label{sec:methods}

\subsection{Dynamics and freeze-out}

Relativistic kick--drift--kick leapfrog with $\dot{\mathbf x}=\mathbf p/E$;
the Hamiltonian is separable so the scheme is symplectic.  Initialization:
periodic cube at density $n$, exact color balance (global singlet),
Maxwell--J\"uttner momenta at $T_0$, Andersen-type equilibration.
Production runs are NVE plus an operator-split Hubble map with scale factor
$a(t)=1+H_0 t$:
\begin{equation}
\mathbf x \to \frac{a_{n+1}}{a_n}\mathbf x,\qquad
\mathbf p \to \frac{a_n}{a_{n+1}}\mathbf p,\qquad
\lambda_D \to \lambda_{D,0}\, a_{n+1},
\label{eq:hubble}
\end{equation}
the comoving screening following from $\lambda_D\sim 1/gT$ with
$T\sim 1/a$.  Cooling outruns dilution, the effective coupling
$\Gamma$ grows, and clustering is \emph{emergent}: the Hubble redshift is
the dissipation channel that permits classical two-body capture.  No
binding rule is imposed.

\subsection{Reactions: Poisson hazards with photon and gluon channels}
\label{sec:reactions}

Matching-color, same-flavor $\qqbar$ pairs within $d_{\rm ann}$ accrue
annihilation hazard (singlet weight $1/3$ folded in):
\begin{equation}
\lambda_{\rm pair} = \frac{1}{3}\times
\begin{cases}
\lambda_{\rm scatt} & E_{\rm pair}>0 \ \text{(in-plasma)},\\[2pt]
1/\tau_{\rm bound} & E_{\rm pair}<0 \ \text{(formed meson or}\\
& \text{\ $\qqbar$ sub-pair in an exotic)},
\end{cases}
\label{eq:hazards}
\end{equation}
where $E_{\rm pair}=M_{\rm inv}-m_i-m_j+V_{ij}$.  Each event branches:
with probability $P_\gamma$ the pair converts to \emph{photons}, whose
four-momentum leaves the system permanently (active at all temperatures ---
the late-time meson-depletion mechanism, against which baryons are
protected by baryon-number conservation); otherwise it converts to
\emph{gluons}, which while $\Teff>\Tchem$ re-inject a matching-color
$\qqbar$ pair carrying the removed four-momentum with flavor drawn from
thermal weights at $\Teff$ (a classical strangeness-enhancement analog),
and below $\Tchem$ are switched off.  All rates are continuous hazards
($P=1-e^{-\lambda\,dt}$), so results are timestep-independent by
construction.  An emergent consequence: internal annihilation converts
tetraquark $\to$ meson and pentaquark $\to$ baryon with lifetime
$\sim 3\tau_{\rm bound}$, giving formation-time versus survival-time ratios
a tunable handle.

\subsection{Species identification}
\label{sec:cluster-id}

A multiset of labels contains a singlet iff
$n_r-n_{\bar r}=n_g-n_{\bar g}=n_b-n_{\bar b}$.  Per late-phase snapshot:
(i) candidate edges between channel-attractive, pairwise Hill-bound pairs
($M_{\rm inv}-m_i-m_j+V_{ij}<0$); (ii) union--find pre-clusters; (iii)
exact minimum-energy partition of each pre-cluster into color-neutral
parts or free singletons, in the spirit of the early-cluster-recognition
algorithm of Ref.~\cite{Dorso:1993zz}; a part is bound iff
$E_{\rm part}<-\Delta E_{\rm th}$.  The tetraquark-versus-two-mesons
ambiguity is resolved by the partition search itself: a $(2,2)$ part
survives only if its four-body energy beats every meson+meson split.  A
persistence filter (identical constituent set in $\ge 3$ consecutive
snapshots) removes transients.  Clusters are tagged by flavor content;
tetraquarks and pentaquarks whose net flavor cannot be carried by
$\qqbar$ or $qqq$ (e.g.\ $cc\bar u\bar d$, the $\Tcc$ analog) are classified
\emph{manifestly exotic} --- the same classification the quantum framework
makes with conserved charges in Sec.~\ref{sec:tier1}.

The formation \emph{pathway} of every (anti)baryon is recorded: the
\emph{diquark-first fraction} --- baryons whose history shows an internal
same-type pair Hill-bound at least $1\fm/c$ before the full cluster
appears --- is the dynamical generalization of the static probability of
Ref.~\cite{Lebed:2016jvq}.
